\documentclass[10pt,conference,compsocconf]{IEEEtran}

\usepackage{hyperref}
\usepackage{graphicx}
\usepackage{amsmath}
\usepackage{amssymb}



\begin{document}
\title{Predicting Heart Disease from Survey Data}

\author{
  Sacha Heizmann, Philipp Glaser, Gaspard Héliot \\
  \textit{École Polytechnique Fédérale de Lausanne, Switzerland}\\
  
}

\maketitle

\begin{abstract}
  We predict coronary heart disease (MICHD) from 321 health and
  lifestyle features of the 2015 BRFSS survey ($328\,135$ training
  respondents, $8.8\%$ positive). We implement the six required methods
  with NumPy only and compare them under one protocol. Because the
  classes are strongly imbalanced we optimize the $F_1$ score and choose
  the decision threshold on held-out data. The only preprocessing step
  that clearly helped was recoding the survey's ``don't know'' and
  ``refused'' answer codes ($+0.013$ $F_1$, better in all five folds);
  squared terms, missing-value indicators and longer training did not
  help beyond noise. Our final regularized logistic regression reaches
  $F_1 = 0.425 \pm 0.006$ in 5-fold cross-validation.
\end{abstract}

% ------------------------------------------------------------------
\section{Introduction}

Cardiovascular disease is a leading cause of death, and early detection
from lifestyle and health information could support prevention. We
use the Behavioral Risk Factor Surveillance System (BRFSS), a US
telephone survey~\cite{brfss}, to predict whether a respondent
reported coronary heart disease or a heart attack. The task is binary
classification with a rare positive class. Our goals were to implement
the basic methods of the course correctly, to build a careful
evaluation protocol, and to test which preprocessing decisions
actually matter, using ablations rather than intuition.

% ------------------------------------------------------------------
\section{Data}
\label{sec:data}

The training set has $N=328\,135$ respondents and $D=321$ features;
the test set has $109\,379$ respondents. Only $8.8\%$ of the training
labels are positive, so a classifier that always answers ``no disease''
reaches $91.2\%$ accuracy while detecting nobody. Accuracy is
therefore a misleading metric, and we use the $F_1$ score, the
harmonic mean of precision and recall of the positive class; the
competition leaderboard also ranks by $F_1$ and additionally shows
accuracy.

The data is incomplete: $44.8\%$ of all values are missing, and $147$
features are missing in more than half of the respondents. Many answers
use special codes: yes/no questions also contain $7$ (``don't know'')
and $9$ (``refused''), and the ``days in the last month'' questions use
$88$ for ``none'' and $77$/$99$ for ``don't know''/``refused''. Read as
numbers, these codes act as large outliers. Some features describe the
survey process (interview date, sequence number, sampling unit), not
the respondent.

% ------------------------------------------------------------------
\section{Models and Methods}
\label{sec:methods}

% Tables are defined here, on page 1: a two-column table (table*) can only
% appear at the top of a page AFTER the one where it is defined.
\begin{table}[htbp]
  \centering
  \footnotesize
  \caption{Required methods on one $80/20$ split with baseline
  preprocessing (no special-code recoding; $173$ columns).
  $\gamma=0.05$ (GD), $10^{-4}$ (SGD, $10^5$ single-sample steps),
  $0.5$ (logistic), $300$ iterations, $\lambda=10^{-4}$.}
  \label{tab:methods}
  \begin{tabular}{|l|c|c|}
    \hline
    Method & Best $F_1$ & Threshold\\
    \hline
    \texttt{mean\_squared\_error\_gd} & 0.406 & 0.209\\
    \texttt{mean\_squared\_error\_sgd} & 0.399 & 0.216\\
    \texttt{least\_squares} & 0.408 & 0.217\\
    \texttt{ridge\_regression} & 0.406 & 0.199\\
    \texttt{logistic\_regression} & 0.412 & 0.201\\
    \texttt{reg\_logistic\_regression} & 0.412 & 0.202\\
    \hline
  \end{tabular}
\end{table}

\begin{table*}[t]
  \centering
  \footnotesize
  \renewcommand{\arraystretch}{1.1}
  \caption{5-fold cross-validation of regularized logistic regression
  ($\lambda=10^{-4}$, $\gamma=0.5$, $300$ iterations, seed $0$), mean $\pm$
  std over folds. $\Delta F_1$: per-fold difference to configuration (c).
  Honest $F_1$: threshold chosen on the other half of the fold.}
  \label{tab:cv}
  \begin{tabular}{|l|c|c|c|c|}
    \hline
    Configuration & Best $F_1$ & Honest $F_1$ & $\Delta F_1$ vs.\ (c)
      & Folds better than (c)\\
    \hline
    (a) baseline & $0.412 \pm 0.005$ & $0.410 \pm 0.006$
      & $-0.0132 \pm 0.0030$ & 0 of 5\\
    (b) + special codes recoded & $0.425 \pm 0.007$ & $0.422 \pm 0.006$
      & $+0.0002 \pm 0.0005$ & 3 of 5\\
    (c) + administrative features removed & $0.425 \pm 0.006$
      & $0.422 \pm 0.006$ & -- & --\\
    (d) (c) + missing-value indicators & $0.426 \pm 0.008$
      & $0.425 \pm 0.008$ & $+0.0007 \pm 0.0018$ & 3 of 5\\
    (e) (c) + squared terms & $0.425 \pm 0.008$ & $0.422 \pm 0.006$
      & $+0.0001 \pm 0.0018$ & 3 of 5\\
    \hline
  \end{tabular}
\end{table*}

\subsection{Implemented methods}

The linear methods minimize
$L(w)=\frac{1}{2N}\sum_n (y_n-x_n^\top w)^2$ by GD, by SGD (one sample
per step) or through the normal equations: least squares solves
$X^\top Xw=X^\top y$, ridge regression ($+\lambda\|w\|^2$) solves
$(X^\top X+2N\lambda I)w=X^\top y$. With the baseline preprocessing
$X^\top X$ is rank-deficient (rank $167$ of $173$), because the survey
contains calculated variables that are exact combinations of other
columns, so least squares uses the pseudo-inverse,
$w=(X^\top X)^{+}X^\top y$. Logistic regression ($y\in\{0,1\}$)
minimizes by gradient descent
\begin{equation}
  L(w)=\frac1N\sum_{n=1}^{N}\Bigl[\ln\bigl(1+e^{x_n^\top w}\bigr)-y_n x_n^\top w\Bigr],
\end{equation}
computed with $\text{logaddexp}$ to avoid overflow; the regularized
version adds $\lambda\|w\|^2$, which keeps the weights finite when the
data is separable (lecture~5b)~\cite{cs433}. The loss returned by the
regularized methods excludes the penalty.

\subsection{Preprocessing}

All statistics are learned on training rows only and then applied
unchanged to validation and test rows, so no information from held-out
data enters the model. (1) \emph{Special codes:} in features whose largest value is $9$, $99$,
$999$ or $9999$, the codes $7$/$9$, $77$/$99$, $777$/$999$, $7777$/$9999$
become missing and $88$/$888$/$8888$ (``none'') become $0$ ($234$ of the
$321$ features). (2) We drop the nine administrative features (state code,
interview date, sequence number, sampling unit, \dots), the features
that are missing in more than $50\%$ of the rows, and constant
features. (3) Remaining missing values are replaced by the median.
(4) Features are standardized, and a constant column is added as bias.
The result has $163$ columns. Labels are converted from $\{-1,1\}$ to
$\{0,1\}$ for training and back for the submission file.

\subsection{Evaluation protocol}

Following lecture~4a we use $M$-fold cross-validation ($M=5$, fixed
seed), fitting the preprocessing inside each fold. Since a model that
outputs probabilities almost never exceeds $0.5$ when only $8.8\%$ of
the labels are positive (at threshold $0.5$ we measured $F_1=0.17$), we
treat the decision threshold as a hyperparameter and choose the one that
maximizes $F_1$. Configurations are compared on the same folds, because
the per-fold \emph{differences} vary much less than the scores
themselves. To estimate how much tuning the threshold on the scored rows
inflates $F_1$, we also report an ``honest'' value: the threshold is
chosen on one half of a validation fold and scored on the other half
(and vice versa). The final threshold is chosen on the pooled
out-of-fold predictions of all $328\,135$ rows, which is more stable than
a per-fold choice.

% ------------------------------------------------------------------
\section{Results}
\label{sec:results}

\emph{Methods.} Table~\ref{tab:methods} compares the six methods. All
six methods reach $F_1$ between $0.399$ and $0.412$; the gaps are
within the noise of a single split (about $0.005$ to $0.008$ from fold
to fold), so we cannot rank them. SGD diverged with step size $10^{-3}$
(weight norm above $10^{9}$); $10^{-4}$ and $10^{-5}$ were stable. We
use regularized logistic regression, which models the label directly
(lecture~5b).

\emph{Ablation.} Table~\ref{tab:cv} compares the preprocessing
configurations in 5-fold cross-validation. Recoding the special codes
raises $F_1$ from $0.412$ to $0.425$ ($0.423$ on a single $80/20$
split), is better in all five folds, and is the only change that helps.
Removing the administrative features, adding indicators for features
with at least $5\%$ missing values ($79$ indicators) and adding squared
terms for $58$ continuous-looking features each change $F_1$ by less
than the fold-to-fold noise, and so does training longer ($1000$ and
$3000$ instead of $300$ iterations give $0.423$ and $0.424$ instead of
$0.424$ on the single split). Our final setting reaches
$F_1=0.425\pm0.006$ in cross-validation ($0.422\pm0.006$ with the
honest threshold estimate), so tuning the threshold on the scored rows
inflates $F_1$ by only $0.001$--$0.003$.

\emph{Final model.} The pooled out-of-fold threshold is $0.195$
(pooled $F_1=0.4250$). On the test set the model predicts $13.9\%$ of the
respondents as positive, almost the same as the $13.9\%$ observed in
cross-validation, so the threshold transfers to the model trained on
all data. On the leaderboard the submission reaches $F_1=0.435$
(accuracy $0.872$), slightly above our cross-validation estimate of
$0.425\pm0.006$.

% ------------------------------------------------------------------
\section{Discussion}

\emph{Strengths.} The pipeline is simple and cheap ($72$ seconds from
the raw files to the submission), every design decision is backed by a
cross-validated comparison, and recoding the special codes is an
interpretable fix for a real property of survey data.

\emph{Why the score plateaus.} Several unrelated changes leave $F_1$
near $0.42$, which suggests that the linear model class is the limit
(high bias, lecture~4b). The six methods are probably so close because
the threshold search compensates for the sensitivity of least squares
to class proportions (lecture~5a); we did not test this separately.

\emph{Limitations.} We did not test one-hot encoding of category-like
features (answer code $3$ is not ``more'' than code $1$), non-linear
models or other losses. The special-code rule is a heuristic that may
misclassify some columns, four survey-design variables (three sampling
weights and a stratum code) remain among the features, and five folds
give only a rough noise level.

% ------------------------------------------------------------------
\section{Summary}

We implemented and compared the six required methods on BRFSS data. The
decisive choices were to evaluate with $F_1$ and a tuned threshold, and
to recode the survey's special answer codes, which improved $F_1$ by
$0.013$ in every fold. Regularized logistic regression on $163$ columns
reaches $F_1=0.425\pm0.006$ in cross-validation and $F_1=0.435$ on the
test set. The ablations show that the remaining ideas we
tried are within noise, pointing to a richer model class or encoding as
the next step. \texttt{run.py} reproduces our submission and
\texttt{cv.py} and \texttt{ablation.py} the tables (see the
\texttt{README}; \url{https://github.com/sachaheizmann/ml-project1}).

\begin{thebibliography}{1}
\bibitem{cs433} N.~Flammarion and M.~Jaggi, \emph{CS-433 Machine
  Learning, lecture notes}, EPFL, 2026.
\bibitem{brfss} Centers for Disease Control and Prevention,
  \emph{Behavioral Risk Factor Surveillance System, 2015 annual data},
  \url{https://www.cdc.gov/brfss/annual_data/annual_2015.html}.
\end{thebibliography}

\end{document}
